\documentclass[11pt,a4paper]{article}
\usepackage[T1]{fontenc}
\usepackage{lmodern,amsmath,amssymb,amsthm,mathtools,booktabs,longtable,microtype,xcolor}
\usepackage[margin=25mm]{geometry}
\usepackage[colorlinks=true,linkcolor=blue!45!black,citecolor=blue!45!black,urlcolor=blue!45!black]{hyperref}
\newtheorem{theorem}{Theorem}[section]
\newtheorem{lemma}[theorem]{Lemma}
\newtheorem{proposition}[theorem]{Proposition}
\newtheorem{corollary}[theorem]{Corollary}
\theoremstyle{definition}\newtheorem{definition}[theorem]{Definition}
\theoremstyle{remark}\newtheorem{remark}[theorem]{Remark}
\newcommand{\Q}{\mathbb Q}\newcommand{\Z}{\mathbb Z}
\newcommand{\HH}{\mathcal H}\newcommand{\ord}{\operatorname{ord}}
\DeclareMathOperator{\Norm}{Norm}
\setlength{\emergencystretch}{2em}
\title{A forced prime condition and an explicit quartic in 392 variables}
\author{}\date{8 November 2026}
\begin{document}\maketitle
\begin{abstract}
We specify an integer quartic in 392 total variables by a complete coefficient
file and a quadratic system of 364 equations. Its rational input projection
is the integers under the arithmetic input contract stated below. The new
argument replaces a three-variable size condition by a single positive
rational parameter with a uniformly large odd valuation at a fixed prime.
The scalar local test then provides both the required lower bound on an
integer index and an independent exclusion of the exceptional auxiliary
elliptic points. Soundness and completeness are proved separately.
The polynomial has 6204 nonzero ordinary monomials, content one, and
coefficient height less than $2^{1791357}$. Exact construction data and a
separate verification program accompany the proof.
\end{abstract}
\tableofcontents

\section{Statement and scope}
For $P\in\Z[X,Y_1,\ldots,Y_m]$, its rational input projection means
\[
 \{x\in\Q:\text{there exist }y_1,\ldots,y_m\in\Q
                  \text{ with }P(x,y_1,\ldots,y_m)=0\}.
\]
A total variable count includes the input $X$. Coefficient height means the
largest absolute value of a coefficient after all fixed integers have been
substituted and equal coordinate monomials have been collected.

\begin{theorem}\label{thm:main}
The files \nolinkurl{quadratic-system392.json} and
\nolinkurl{polynomial392.json} specify 364 integer quadratic polynomials
$Q_j$ in 392 variables, and the polynomial $P=\sum_jQ_j^2$.
The degree of $P$ is four, its content is one, and it has 6204 nonzero
coordinate monomials. Its largest coefficient has 1791357 binary digits.
Under the arithmetic input contract in Section~\ref{sec:inputs}, its
rational input projection is $\Z$.
\end{theorem}
The coefficient assertions are exact finite algebraic statements. The
projection assertion is a conditional theorem: the computational programs
do not establish the arithmetic hypotheses. The changes made here do not
assert a bijection with all witness assignments of the starting system.
They prove a sound integer extraction and construct witnesses for every
integer input.

\section{Arithmetic inputs and the retained system}\label{sec:inputs}
Put
\[
 K=\Q(\rho),\quad \rho^2=2,\quad d_0=3-\rho,\quad D=d_0^2,
 \qquad E:y^2=x^3-Dx.
\]
The fixed points are
\[
 R_0=(1+\rho,2\rho),\qquad
 P=2R_0=\left(\frac{57-28\rho}{8},\frac{-660+417\rho}{32}\right).
\]
For a point in its domain, write $f=1/x$ and $z=-x/y$, extended by
$f(O)=z(O)=0$ near the identity. Heights are absolute logarithmic Weil
heights. At an unramified place above $p$, the valuation is normalized by
$v(p)=1$.

The arithmetic input contract consists of the following assertions,
with the original sources and applications identified in
\cite[Sections \nolinkurl{subsec:v12-inputs} and
\nolinkurl{sec:v13}]{GlazerV13}.
\begin{enumerate}
\item The equality $2E(K)=\Z P$ holds, and
$h(x([n]P))\ge n^2-6$ for every nonzero integer $n$.
\item For $t=u/v$ in lowest terms, $v>0$, outside
$\{0,1,2401,-343/32\}$, let $M_5(t)$ be the product of primes
$p>76176$ for which a positive odd valuation occurs in at least one of
\[
 v,\quad u,\quad u-v,\quad u-2401v,\quad32u+343v.
\]
The contact-height inequality is
\[
 h(t)\le\HH M_5(t)^{19},\qquad
 \HH=2^{110}210^{1800}\left(\prod_{p\le76176}p\right)^{19}+128.
\]
The exact numerical bound used here is $\HH<2^{2093448}$.
\item The pole equations in the source, with input $b$ and $B=17b$,
\[
 U=b\frac{1-bt^2}{1+bt^2},\quad h_0=e(B-U^2),\quad
 p^2=B(1-e)(1+3e),\quad
 Bq^2=(s^2-Bh_0)(s^2+3Bh_0),
\]
exclude negative odd $v_\ell(b)$ when $\ell\notin\{2,3,17\}$.
They admit the specialization $t=0$ if $0<b<17$, the reduced denominator
of $b$ is a square, $b$ is a square in each of $\Q_2,\Q_3,\Q_{17}$,
and its valuations there are at least $3,1,2$, respectively. We use
the exact normalized seven-coordinate form of these equations in the
specified starting system.
\item The source simultaneous prime-pattern and conjugate reduction-order
conditions can be realized with arbitrary fixed odd avoidance modulus.
In the notation of that source, selected primes are $1$ modulo
$8\cdot3\cdot17$, the orders satisfy
$v_2(d_-)=6$, $v_2(d_+)=2$ and
$\gcd(d_-(p),d_+(p'))=4$. The required fixed congruences and real
approximation conditions are simultaneous.
\end{enumerate}
As in the quartic source, we retain the following quantitative form of the
last input for the separation step. For the fixed distinct coefficients
$c_j$, the admissible patterns $p_j=u+c_jw^2$ with
$N/2\le u\le N$ and $1\le w\le\tau(N)\sqrt N$ have total weight at least
$c_0N^{3/2}\tau(N)$, individual weights are
$O((\log N)^k)$, and $\tau(N)$ may decrease arbitrarily slowly after
all fixed congruences and avoidance conditions are imposed. This is an
arithmetic premise; finite coefficient checks do not prove it. The precise
source locations are Proposition \nolinkurl{prop:v5-artin-patterns}, the
final diagonal argument in \nolinkurl{subsec:v5-diagonal}, and the
separation application in the retained Version 9 portion of
\cite{GlazerV13}.

\subsection{The unchanged comparison interface}
The starting quadratic system is the 394-variable system of
\cite{GlazerV13}, at the commit specified in Section~\ref{sec:data}.
We retain all its point charts, comparison equations, three height tests,
and their normalized poles. The following are the properties of that
unchanged part used here. Their proofs and precise local applications
are given in the Version 13 sections on reciprocal charts, the enlarged
$H$ domain, scaled charts, and the two final height contradictions.

The four root points are
\[
 Q_1=[n]P,\quad S=Q_2=[a]P,\quad Q_3=[t]P,\quad Q_4=[b]P,
 \qquad n,a,t,b\text{ nonzero odd integers}.
\]
The guarded reciprocal subtraction chart gives
\[
 G=S-P=[s]P,\qquad s=a-1\ne0.
\]
The auxiliary finite point is $H=[m]R_0+T$ with $m\in\Z$ and
$T\in E[2]$. Its retained chart is
\[
 x(H)f(H)=1,\qquad H=(x,xS_H),\qquad S_H^2=x-Df(H).
\]
It follows immediately, before assigning meaning to any ratio slot, that
$f(H)\ne0$ and $q=\Norm_{K/\Q}(f(H))\ne0$.

After excluding
\begin{equation}\label{eq:exceptional}
 E[2]\ \cup\ (P+E[2])\ \cup\ (-P+E[2]),
\end{equation}
the retained ratio equations identify their scalar slots with
\[
                  \frac{z(G)}{z(H)}=z_0+z_1\rho.
\]
The proof also identifies the primary and alternate comparison slots,
including the identity branches of their numerators. It uses the integral
group identities
\[
 [m]U_i=[2u_i]H,\qquad
 [m-2]V_i=[2(u_i-n_i)](H-P),
\]
so that $h=m/2$ is only a rational index; no fractional group multiple is
asserted. The resulting order-forty tests make the three contact radicals
divide
\begin{equation}\label{eq:F}
 F=t^2-a(a-2)tb+(a-2)^2b^2-1
\end{equation}
to order forty. The order-ten tests make them divide the nonzero reduced
numerator of $n-\beta$ to order ten, if that difference is nonzero.
The common height consequence is
\begin{equation}\label{eq:height}
 B_*^2-6<7\HH |N|^{19/k},\qquad
 B_* =\max(|n|,|a|,|t|,|b|),
\end{equation}
whenever their $k$th powers divide $N\ne0$. The three scales remain
$1,C_0^2,(1+C_1^2)^2$ for
$\prod_i x(Q_i)=C_0+C_1\rho$.
None of these equations, scales, or comparison orders is changed below.

For completeness the source construction permits any fixed finite list
of formal-depth conditions on $H$ and $G$, imposed before the later
contact primes are chosen. Let $\sigma$ be conjugation on $K$, and put
\[
 G_{\mathrm{real}}
 =\overline{\{([n]P,([n]P)^\sigma):n\in\Z\}}
 \subset E(\mathbb R)\times E^\sigma(\mathbb R).
\]
Write $G_{\mathrm{real}}^0$ for its identity component and
$m_0=|G_{\mathrm{real}}/G_{\mathrm{real}}^0|$.
A progression with residue and modulus divisible by $m_0$ has dense
positive tail in $G_{\mathrm{real}}^0$.
Thus the prescribed finite conditions and small real values
of $f(H)$ and $z(G)/z(H)$ can hold at the same global points. The fixed
modulus can include any further finite odd support. This simultaneous
version, rather than separate choices at different places, is used in
Section~\ref{sec:complete}.

\section{A positive parameter with a forced odd valuation}
Define
\begin{align}
 \beta(X)&=51(2X+1)^2+2504,\label{eq:beta}\\
 r(X)&=\beta(\beta-1)(\beta-2401)(32\beta+343)(2\beta-1).\label{eq:r}
\end{align}
For every rational $X$, all five factors are positive. In particular
$r(X)>0$ without any integrality assumption.

\begin{lemma}\label{lem:prime}
The integer $\ell=80989$ is prime, is greater than $76176$, and is
$1$ modulo $17$. For every $X\in\Q$,
\[
 v_\ell(r(X))=
 \begin{cases}0,&v_\ell(X)\ge0,\\10v_\ell(X),&v_\ell(X)<0.
 \end{cases}
\]
\end{lemma}
\begin{proof}
Trial division through $\lfloor\sqrt{80989}\rfloor=284$ certifies
primality. For each of the five contact roots $\theta$, Euler's criterion gives
\[
 \left(\frac{(\theta-2504)/51}{\ell}\right)=-1.
\]
The exact residues and modular powers are listed in the verification
file and independently recomputed by the checker. Hence, if $X$ is
$\ell$-integral, $\beta(X)$ misses all five roots of the product
modulo $\ell$, and $r$ is a unit. If $v_\ell(X)<0$, the leading term
of $\beta$ gives $v_\ell(\beta)=2v_\ell(X)$. Every contact factor
then has that valuation, proving the second assertion.
\end{proof}

\begin{lemma}\label{lem:normalizer}
Let $N\ge1$ be an integer and put
\[
 A=\ell^{2N+1},\qquad
 R=\frac{A r}{1+\ell^2 r},\qquad c=R+4896.
\]
For every rational $X$, $R,c>0$ and
\begin{equation}\label{eq:forcednu}
 v_\ell(R)\in\{2N+1,2N-1\}.
\end{equation}
Every prime occurring in $M_5(\beta)$ is different from $\ell$ and
has positive odd valuation on $R$; at such a prime $c$ is a unit.
For an integer $X$,
\begin{equation}\label{eq:smalllocal}
 v_2(R)=2,\quad v_3(R)=1,\quad v_{17}(R)=0,
 \qquad c/R\in\Q_p^{\times2}\quad(p=2,3,17).
\end{equation}
\end{lemma}
\begin{proof}
The denominator $1+\ell^2r$ is positive. If $v_\ell(r)=0$, it is a
unit. Otherwise $v_\ell(r)\le-10$, so its valuation is
$2+v_\ell(r)$. These two cases give \eqref{eq:forcednu}.

The reduced denominator of $\beta$ has even valuation at every prime
larger than $76176$. Therefore its contact radical only consumes positive
odd valuations at the four finite roots defining $M_5$. Distinct roots
of the five factors in \eqref{eq:r} cannot coincide modulo a prime
larger than $76176$: the relevant nonzero differences have no such
prime divisor. At a consumed prime, consequently, $v_p(r)>0$ is odd.
Lemma~\ref{lem:prime} excludes $p=\ell$. Both $A$ and
$1+\ell^2r$ are then units, so $v_p(R)=v_p(r)$ and $c$ is a unit,
since every prime factor of $4896$ is $2$, $3$, or $17$.

For integral $X$, the odd integer $\tau=2X+1$ gives
$v_2(r)=2$, $v_3(r)=1$ and $r\equiv15\pmod{17}$. These follow directly
by reducing the five displayed factors modulo $8,9,17$, respectively.
The new denominator is a unit at $2,3$, and is $16$ modulo $17$.
The numerator multiplier is a unit at all three primes. Thus the
valuations of $R$ are as stated. Finally
\[
 v_2(4896/R)=3,\qquad v_3(4896/R)=1,\qquad
 v_{17}(4896/R)=1.
\]
The unit $1+4896/R$ is therefore a local square: it is $1$ modulo $8$
at two and $1$ modulo the maximal ideal at the two odd primes.
\end{proof}

\section{Local extraction before interpreting the ratio slots}
The retained scalar norm equations, with their parameter changed from
$r$ to $R$, are
\begin{equation}\label{eq:normtree}
 B_s=\frac{z_0^2+Rz_1^2}{R^{32}},\qquad
 W=\frac{B_s^2+Rq^2}{R^{32}},\qquad
 \mathfrak b=(R+4896)RW^2.
\end{equation}
Initially $z_0,z_1$ in this display mean the stored rational slots. Their
ratio interpretation is not used in the next lemma. Since $q\ne0$ and
$R>0$, we have $W>0$ and $\mathfrak b>0$. Hence the normalized pole
has a nonzero input before any logarithm denominator is interpreted.

\begin{lemma}\label{lem:depth}
Suppose $p>76176$, $\nu=v_p(R)>0$ is odd, and the pole equations
for \eqref{eq:normtree} hold. Then
\[
 v_p(B_s)>\frac{63\nu}{4},\quad
 v_p(q)>\frac{61\nu}{4},\quad
 v_p(z_0)>\frac{191\nu}{8},\quad
 v_p(z_1)>\frac{187\nu}{8}.
\]
\end{lemma}
\begin{proof}
The correction $R+4896$ is a unit. The valuation of $\mathfrak b$ is
$\nu+2v_p(W)$, an odd integer. The pole obstruction excludes negative
values, so $v_p(W)>-\nu/2$. The summands in each of the two norm
numerators have valuations of opposite parity. They cannot cancel;
a zero summand is assigned valuation $+\infty$. The second equation
in \eqref{eq:normtree} therefore gives
\[
 2v_p(B_s)>\frac{63\nu}{2},\qquad
 \nu+2v_p(q)>\frac{63\nu}{2}.
\]
Apply the same argument to the first numerator to obtain
$2v_p(z_0)>191\nu/4$ and
$\nu+2v_p(z_1)>191\nu/4$. These are the asserted bounds.
\end{proof}

\begin{lemma}\label{lem:exclude}
The equations with the forced parameter exclude every exceptional
point in \eqref{eq:exceptional}, before the ratio slots need to be
interpreted.
\end{lemma}
\begin{proof}
The points $O$ and $(0,0)$ are already excluded by $xf(H)=1$.
For every remaining exceptional point, its abscissa is among
\[
 \pm d_0,\quad c_P,\quad -D/c_P,\quad
 d_0\frac{c_P+d_0}{c_P-d_0},\quad
 -d_0\frac{c_P-d_0}{c_P+d_0},
 \qquad c_P=\frac{57-28\rho}{8}.
\]
These are the usual translations by the three two-torsion points.
The factors and their inverses are integral outside
$\{2,3,7,17,41\}$, because
\[
 \Norm(d_0)=7,\quad \Norm(c_P)=\frac{41^2}{64},\quad
 \Norm(c_P-d_0)=\frac{17^2}{64},\quad
 \Norm(c_P+d_0)=\frac{63^2}{64}.
\]
Thus $f(H)$ is a unit at every place above $\ell$, and $q$ is an
$\ell$-unit. This contradicts Lemma~\ref{lem:depth} and
\eqref{eq:forcednu}. No meaning of $z_0,z_1$ was used.
\end{proof}

\begin{proposition}\label{prop:index}
Every solution has a nonzero even integer $s=a-1$ satisfying
\[
 v_\ell(s)\ge
 L_N:=\left\lfloor\frac{187(2N-1)}8\right\rfloor+1.
\]
Moreover $M_5(\beta)^{24}$ divides $s$.
\end{proposition}
\begin{proof}
The retained reciprocal guard gives $G=[s]P$ with $s\ne0$ even.
Lemma~\ref{lem:exclude} now permits the ratio interpretation
$z_0+z_1\rho=z(G)/z(H)$. At $\ell$, positive valuation of $q$
supplies a place at which $v(f(H))>0$, so $H$ is formal there.
The rational coefficient depths from Lemma~\ref{lem:depth} give
\[
 v(z(G))\ge v(z(H))+L_N.
\]
If $e=v(x(G))$, the identity
$z(G)^2=x(G)/(x(G)^2-D)$ shows that $e\ne0$.
The duplication formula
\[
 x(2G)=\frac{(x(G)^2+D)^2}{4x(G)(x(G)^2-D)}
\]
gives $v(x(2G))=-2v(z(G))$ for both signs of $e$. Hence $2G$ is
formal and $v(z(2G))=v(z(G))$. Apply the formal logarithm, an isometry
at this positive depth, to the integral group identity
\[
                       [m](2G)=[4s]H.
\]
The excluded set implies $m\ne0$. Since $m/2$ is $\ell$-integral,
we obtain
$v_\ell(s/(m/2))\ge L_N$ and consequently $v_\ell(s)\ge L_N$.

For a prime in $M_5(\beta)$ the same proof applies, by
Lemma~\ref{lem:normalizer}. Here $\nu\ge1$ is odd, so both coefficient
depths are at least $24$. The formal comparison gives $v_p(s)\ge24$
at each such prime, which proves the divisibility assertion.
\end{proof}

\section{The numerical choice and integer extraction}
We now fix
\[
              N=27464,\qquad A=80989^{54929}.
\]
Then $2N-1=54927$ and $L_N=1283919$. The entirely integral
certificates
\[
 80989^{4096}>2^{66787},\qquad
 66787\cdot1283919-4096\cdot20934550=1181453>0
\]
show that every nonzero index in Proposition~\ref{prop:index} satisfies
\begin{equation}\label{eq:large}
 |s|\ge80989^{1283919}>2^{20934550}>(80\HH)^{10}.
\end{equation}
Since $|a|=|s+1|\ge |s|-1$, this implies
\begin{equation}\label{eq:alarge}
 |a|>(40\HH)^{10},\qquad
 \log(48\HH)<|a|/40.
\end{equation}
The scalar contact-height inequality and the divisibility just proved give
\begin{equation}\label{eq:betasmall}
 h(\beta)\le\HH |s|^{19/24}<|s|/80\le |a|/40.
\end{equation}
The middle inequality follows from
$|s|^{5/24}>(80\HH)^{50/24}>80\HH$.

For clarity, the elementary growth fact needed for the integer extraction
is proved here, rather than included among the arithmetic inputs.

\begin{lemma}\label{lem:lucas}
If integers $a,t,b$ satisfy $|a|>7$, $b\ne0$, and
$t^2-a(a-2)tb+(a-2)^2b^2=1$, then $\log|t|>|a|-3$.
\end{lemma}
\begin{proof}
Set $A_0=|a|$ and $c=\operatorname{sgn}(a)(a-2)b$. Then
$t^2-A_0tc+c^2=1$. The nonzero integers $t,c$ have the same sign,
and $T=|t|$, $C=|c|$ are positive solutions. Define
$U_0=0$, $U_1=1$, $U_{j+1}=A_0U_j-U_{j-1}$.
All positive solutions are unordered adjacent terms of this sequence:
for $M\ge m>0$, the replacement
\[
 (M,m)\longmapsto(m,A_0m-M),\qquad
 A_0m-M=(m^2-1)/M\in[0,m),
\]
preserves the equation and descends to $(1,0)$.
Thus $C=U_j$ and $T\ge U_{j-1}$ for some $j\ge1$.
If $d=|a-2|$, reduction of the recurrence modulo $d$ gives
$U_j\equiv j$ or $(-1)^{j-1}j$, according to the sign of $a$.
Since $d\mid C$, we have $j\ge d\ge A_0-2$.
The recurrence gives $U_r>(A_0-1)^{r-1}$ for $r\ge2$. Hence
\[
 \log T>(j-2)\log(A_0-1)
       \ge(A_0-4)\log7>A_0-3\qquad(A_0\ge8).
\]
The final inequality holds at $8$ and its difference increases.
\end{proof}

\begin{proposition}[Soundness]\label{prop:sound}
Every rational solution of the new quadratic system has integer input.
\end{proposition}
\begin{proof}
All the unchanged comparison slots have their required meaning by
Lemma~\ref{lem:exclude}. If the integer $F$ in \eqref{eq:F} were
nonzero, its order-forty divisibility and \eqref{eq:height} would give
\[
 B_*^2-6<7\HH(5B_*^4)^{19/40}.
\]
As $B_*>4$, this implies $B_*^{1/10}<40\HH$, contrary to
$B_*\ge|a|$ and \eqref{eq:alarge}. Thus $F=0$, and
Lemma~\ref{lem:lucas} yields $\log B_*>|a|-3$.

If $n-\beta\ne0$, let $N_0$ be its nonzero reduced numerator.
The source comparison gives order-ten divisibility and
$|N_0|\le2e^{h(\beta)}B_*$. Its height consequence is
\[
 \log B_*<10\log(48\HH)+19h(\beta)<29|a|/40<|a|-3,
\]
where \eqref{eq:alarge}, \eqref{eq:betasmall}, and the enormous lower
bound on $|a|$ justify the last two inequalities. This is impossible.
Hence $n=\beta$.

Finally, write $2X+1=u/v$ in lowest terms with $v>0$. Integrality of
$51u^2/v^2+2504$ implies $v^2\mid51$, so $v=1$ because $51$ is
squarefree. Since $n=\beta$ is odd, $u$ is odd. Therefore
$X=(u-1)/2\in\Z$.
\end{proof}

\section{Simultaneous completeness}\label{sec:complete}
Fix an integer input $X$. The numbers $R,c$ are now fixed positive
rationals satisfying \eqref{eq:smalllocal}. Let $S$ contain
$2,3,17$ and every prime with nonzero valuation on $cR$ or $R$.
This is a finite set, fixed before the later prime patterns are chosen.

\begin{lemma}\label{lem:prepare}
The same global auxiliary points can be chosen to make
\[
        0<\mathfrak b<17,
\]
while $\mathfrak b$ has square reduced denominator, is a local square
at $2,3,17$, and has valuations there at least $3,1,2$.
These requirements are compatible with every retained comparison and
prime-selection condition.
\end{lemma}
\begin{proof}
At a fixed finite place all coefficients in \eqref{eq:normtree} are
fixed. If the two rational coefficients of $f(H)$ and those of
$z(G)/z(H)$ tend to zero, then $q,B_s,W$ tend to zero. In valuation
terms, if both pairs have depth at least $t$ and $\alpha=v_p(R)$,
then
\[
 v_p(q)\ge2t,\qquad
 v_p(B_s)\ge2t+\min(0,\alpha)-32\alpha,
\]
and
\[
 v_p(W)\ge
 \min\{2v_p(B_s),\alpha+2v_p(q)\}-32\alpha.
\]
These lower bounds tend to infinity with $t$. Choose sufficiently
large depths at every prime in $S$ that $\mathfrak b$ is integral
there and has the three required special lower bounds. Away from
$S$ its valuation is $2v_p(W)$, which is even. Thus the reduced
denominator of $\mathfrak b$ is a square. Its local squareclasses
at $2,3,17$ equal those of $cR$, which are squares by
\eqref{eq:smalllocal}.

These local conditions must be achieved together with the real
condition. In the paired real orbit closure, first choose a positive
index $h_0$ divisible by $64m_0$, in a sufficiently divisible
progression, so that
$H_0=[h_0]P$ has all selected formal depths and is close to both real
identities. Then $q(H_0)$ is arbitrarily small. Fix this point.
Next choose a positive index $s$ divisible by $6$ in a sufficiently
divisible progression, arbitrarily large and with $G=[s]P$ close to both real
identities. As $z(H_0)$ is nonzero, the two embeddings of
$z(G)/z(H_0)$, and therefore its two rational coefficients, are
arbitrarily small. The displayed norm tree is continuous and tends
to zero at these values. Choose the points so that
$0<\mathfrak b(H_0,G)<17/4$ and all finite conditions hold.
Strict positivity follows from $q(H_0)\ne0$.

Keep $G$ fixed. Real continuity supplies a neighborhood of $H_0$
with $0<\mathfrak b<17$, and finite continuity supplies a fixed
congruence $h\equiv h_0\pmod L$ preserving all selected depths.
Choose $64m_0\mid L$. The remaining root labels and comparison
scales are fixed before the new contact primes. A smaller nonempty
real open set with compact closure avoids all comparison poles;
their bounds on that closure are consequently fixed as well.

The retained prime-selection and separation argument allows this
additional fixed modulus. Put
\[
 \Pi=\prod_p p,\qquad
 M_\pm=\operatorname{lcm}_p(d_\pm(p)/4),
\]
where the products and least common multiples run over the selected
primes. The order and avoidance conditions give
$\gcd(M_-,M_+)=1$, $\gcd(M_-,L)=16$, and
$\gcd(M_+,L)=1$. The separation step gives
$\gcd(\Pi,M_+)=1$, and the selected primes avoid $L$.
Since $64\mid h_0$, the following congruences are compatible:
\[
 h\equiv0\pmod{\Pi^{19}M_-},\qquad
 h\equiv1\pmod{M_+},\qquad h\equiv h_0\pmod L.
\]
The resulting progression has residue and modulus divisible by $m_0$,
so its positive tail meets the selected real open set. A single such
final $H$ satisfies the finite depths, the real bound, all comparison
domains, and the contact requirements. This proves simultaneity.
\end{proof}

\begin{proposition}[Completeness]\label{prop:complete}
Every integer input admits a rational solution of the new system.
\end{proposition}
\begin{proof}
Use the unchanged source construction of the roots, charts, and
comparison points with the finite conditions of Lemma~\ref{lem:prepare}.
The pole-realization premise supplies fresh scalar pole witnesses
with preparation specialized to zero. Its exact normalized chart
permits zero private witnesses and divides only by the nonzero
semantic pole input. Thus all five normalized pole equations hold.
The guard on the reciprocal scalar point is also available. Indeed,
write $s=2\cdot3^ku$ with $3\nmid u$. Direct duplication of the
specified point $P$ gives
\[
                z(2P)/3\equiv2+\rho\pmod3.
\]
The integral odd formal multiplication series has linear term $nz$
and terms of degree at least three thereafter. At positive depth,
iteration of multiplication by $3$ and then by $u$ therefore gives
\[
 z([s]P)/3^{k+1}\equiv u(2+\rho)\pmod3.
\]
Using $f=z^2(1-Df^2)$ and $(2+\rho)^2\equiv\rho\pmod3$,
we find that the $\rho$-coefficient of $f([s]P)$ has valuation exactly
$2+2k$. If $f([s]P)=g^2$, this coefficient is $2g_0g_1$.
It is nonzero, so $g_0^{-1}$ exists. This argument also covers negative
$u$ and needs no deleted anchor condition.

Assign the new coordinate its uniquely determined value
$R=Ar/(1+\ell^2r)$, and compute its five successive squares.
The scalar equations hold by definition. The deleted three private
anchor witnesses have no other occurrences. All remaining retained
witnesses are supplied by the same simultaneous construction.
\end{proof}
Because rational squares are nonnegative, $\sum_j Q_j^2=0$ over
$\Q$ is equivalent to all the equations $Q_j=0$. Propositions
\ref{prop:sound} and \ref{prop:complete} therefore prove the
projection part of Theorem~\ref{thm:main}.

\clearpage
\section{Exact quadratic implementation}\label{sec:data}
The starting file is the original-byte
\nolinkurl{candidate394-ir.json} at upstream commit
\begin{center}\small
\nolinkurl{5bbd74f6a479776a332cbe036809d5dddae63720}.
\end{center}
Its SHA-256 digest is
\begin{center}\small
\nolinkurl{7eb02f9913b55d60a790a287616f1bb80d9f93523e9f11e12481c9ab955746a5}.
\end{center}
The complete starting coefficient table is separately supplied and is
independently reconstructed from this system before the new computation.

We use the source coordinate numbers $v_i$ only to make the exact
changes unambiguous; the final files renumber the coordinates from
$1$ through $392$, with the input at coordinate $1$. The raw contact
in the retained source gates is the affine expression
\begin{align*}
 r={}&-105275102108-8412883296v_0-8412883296v_{248}\\
     &-21050600v_{487}-7629700v_{488}-45771136v_{489}.
\end{align*}
The source equations numbered $200,453,454,455$ imply that this
expression equals \eqref{eq:r}. The independent checker proves the
identity as an exact univariate polynomial of degree ten.

Delete the source coordinates $55,56,57$ and their sole equation,
numbered $504$. Add one coordinate $v_{702}=R$ and the quadratic equation
\begin{equation}\label{eq:newgate}
                 R+\ell^2rR-Ar=0.
\end{equation}
It is quadratic because the displayed source expression for $r$ is
affine in retained coordinates. The denominator is always positive,
so \eqref{eq:newgate} has the unique intended rational solution.
Replace exactly the following four source rows:
\begin{center}
\begin{tabular}{cl}
\toprule Source row & New equation \\\midrule
$560$ & $v_{600}-R^2=0$,\\
$566$ & $v_{604}v_{540}-v_{508}^2-Rv_{539}=0$,\\
$512$ & $v_{604}v_{542}-v_{540}^2-Rv_{541}=0$,\\
$492$ & $(v_{600}+4896R)v_{525}+82479248v_{526}=0$.\\
\bottomrule
\end{tabular}
\end{center}
The retained power rows give
$v_{600}=R^2$, $v_{601}=R^4$, $v_{602}=R^8$,
$v_{603}=R^{16}$ and $v_{604}=R^{32}$.
The retained row $565$ gives $v_{539}=v_{509}^2$, and row $511$
gives $v_{541}=q^2$ through the exact identity
\[
 q^2=(f_0^2+2f_1^2)^2-8(f_0f_1)^2.
\]
If $J$ is the positive integer coefficient in row $491$,
$Jv_{525}-v_{542}^2=0$, then the actual semantic pole input is
$-82479248Jv_{526}=(R+4896)RW^2$.
This sign and fixed scale are checked from the literal rows.

The guard $v_{543}v_{549}=1$ is retained. In particular, the
nonzero scalar index has not been replaced by an implicit assumption.
Every other coordinate and residual is unchanged. The final count is
\[
               394-3+1=392,\qquad364-1+1=364.
\]
For each rational-coefficient row, clear denominators and divide by
the gcd of its integer coefficients. Square these primitive rows and
add. The resulting polynomial is the supplied $P$. Its symbolic
coefficient table uses $A=80989^{54929}$ as a fixed integer, never as
an extra variable.

\section{Independent computational verification}
The producer \nolinkurl{build_quartic392.py} constructs the complete
system and polynomial. The separate program
\nolinkurl{verify_quartic392.py} imports no producer or earlier checker.
It reconstructs the upstream 394-variable sum of squares and matches
its public coefficient table. It then checks the full changed system,
including every unchanged row, all coordinate incidences, the affine
raw contact, the normalization and inverse guard. Finally it expands
all 364 squares and compares every coefficient.

The new table has 6225 terms when powers of the fixed parameter are
listed separately. After those powers are substituted and equal
coordinate monomials are collected, it has 6204 nonzero monomials.
These two counts are intentionally distinct. A parameter-independent
coefficient equal to one certifies content one; a nonzero quartic
monomial certifies the exact degree.

The maximum coefficient is the constant coefficient. If $C_{\mathrm{const}}$ is the
sum of the squared constant terms of the 359 unchanged primitive
residuals, then it is exactly
\[
       11082847123849826043664\cdot80989^{109858}+C_{\mathrm{const}},
       \qquad 0<C_{\mathrm{const}}<2^{1514}.
\]
The full integer $C_{\mathrm{const}}$ is included in
\nolinkurl{construction-report392.json}; its definition here also
specifies it without abbreviating decimal digits. The independent
checker compares all specialized coefficients as exact integers and
proves that the maximum has 1791357 binary digits.

The same program checks trial division for $80989$, all five Euler
criterion calculations, the local input residues, and the integral
height margin. Its receipt is
\nolinkurl{independent-verification392.json}.
Run from the construction directory:
\begin{verbatim}
python build_quartic392.py
python verify_quartic392.py
\end{verbatim}
These checks establish the finite algebra and numerical certificates.
The arithmetic hypotheses and the simultaneous construction remain
the mathematical inputs and arguments stated in this paper.

\begin{thebibliography}{9}
\bibitem{GlazerV13}
E.~Glazer, \emph{An explicit quartic defining the integers in the
rationals}, Version 13, repository manuscript and coefficient data,
commit \nolinkurl{5bbd74f6a479776a332cbe036809d5dddae63720},
\url{https://github.com/elliotglazer/diophantine-z-q}.
The Version 13 addition and the separate mathematical composition
reviews identify the enlarged auxiliary-point domain, scalar norm
calculation, pole normalization, and simultaneous completeness
argument used here.
\bibitem{Silverman}
J.~H.~Silverman, \emph{The Arithmetic of Elliptic Curves}, second
edition, Graduate Texts in Mathematics 106, Springer, 2009,
Chapter IV, especially Theorem 6.4.
\end{thebibliography}
\end{document}
